Without a wire, the three identical bulbs in series share the battery voltage $\Delta V$ equally: each gets $\frac{1}{3}\,\Delta V$.

\begin{abcliste}
\abc The wire has no resistance, so Q and R are at the same potential: there is no voltage across B, and the whole current takes the wire. B goes \result{\text{off}}. Only A and C share the battery voltage now, each with $\frac{1}{2}\,\Delta V$ instead of $\frac{1}{3}\,\Delta V$. With more voltage they carry more current, so A and C get \result{\text{brighter}} (and are equally bright).

\abc Now there is no voltage across A and B: both are \result{\text{off}}. C gets the whole battery voltage $\Delta V$: it is \result{\text{brighter}} than without the wire ($\frac{1}{3}\,\Delta V$) and brighter still than in (a) ($\frac{1}{2}\,\Delta V$). It is as bright as one bulb on its own on the battery.
\end{abcliste}

A typical mistake is to think that a bridged bulb still lights because the current has to pass through it. But the bulb has no voltage across it, and the current takes the wire.
